\subsection{Zielgruppenanalyse}\label{subsec:zielgruppenanalyse}

\subsubsection{Nutzersegmente}\label{sec:nutzersegmente}

Die Nutzer:innen der DocSecBox lassen sich in zwei Hauptsegmente unterteilen, die sich in Anforderungen, Nutzungskontexten und technischen Voraussetzungen deutlich unterscheiden.

Das erste Segment umfasst professionelle Anwender:innen in regulierten Branchen -- Ärzt:innen, Anwält:innen, Steuerberater:innen und Notar:innen --, die regelmäßig hochsensible Daten wie Befunde oder Verträge behandeln. Daraus ergeben sich hohe Anforderungen an Nachvollziehbarkeit, Datenschutz und die Einhaltung der DSGVO. Die DocSecBox begegnet diesen Anforderungen mit lückenlosem Audit-Log, klarer Trennung der Zugriffsarten (registriert vs. anonym) und vollständiger Datei-Verschlüsselung. Dieses Segment bildet den primären Fokus des Redesigns.

Das zweite Segment umfasst Mitarbeiter:innen und Führungskräfte, die Dokumente intern und an externe Partner weitergeben -- etwa Prozessbeschreibungen, Schulungsmaterial oder Vertragsentwürfe. Im Fokus stehen effiziente Verteilung, klare Ordnerstrukturen und geringe Zugangshürden für die Empfänger:innen. Die DocSecBox ermöglicht öffentliche und nicht-öffentliche Freigaben.

\subsubsection{Zugriffsarten und Nutzungsprofile}\label{sec:zugriffsarten}

Unabhängig von der Segmentierung unterscheidet die DocSecBox zwei technische Zugriffsarten. Registrierte Nutzer:innen sind authentifiziert und können je nach Berechtigung Ordner erstellen, verwalten und freigeben. Innerhalb dieser Gruppe lassen sich drei Nutzungsprofile unterscheiden: \textit{Nur-Empfänger:innen} (konsumieren ausschließlich Empfangenes), \textit{Nur-Provider:innen} (laden primär hoch) und \textit{Provider und Empfänger} (Zwei-Wege-Nutzung, z. B. Arzt-zu-Arzt-Kommunikation). Anonyme Nutzer:innen empfangen Dateien ausschließlich über freigegebene Links ohne Account. Diese Profile prägen die drei Personas in \autoref{sec:personas}.
